
\documentclass{article}
\usepackage[utf8]{inputenc}
\usepackage{amsmath}
\usepackage{amsfonts}
\usepackage{amssymb}
\usepackage{graphicx}
\usepackage{hyperref}

\title{Enhancing Faithful Confidence Expression in Large Reasoning Models Through Advanced Prompt Engineering}
\author{AI Research Assistant}
\date{June 2024}

\begin{document}

\maketitle

\begin{abstract}
The trustworthiness of Large Language Models (LLMs), particularly Large Reasoning Models (LRMs), in high-stakes applications hinges on their ability to communicate uncertainty reliably. This paper synthesizes findings from recent research, specifically Gani et al. (2026), to explore the concept of Faithful Calibration (FC) � the alignment between an LRM's intrinsic confidence and its linguistically expressed certainty. We discuss the challenges in quantifying FC in complex, multi-step reasoning traces and review the methodological framework proposed by Gani et al., which utilizes multiple intrinsic confidence estimators and a novel prefix-conditioned sampling approach. Key findings indicate that current LRMs struggle with FC, reasoning-specific training can degrade faithfulness, and conventional prompt engineering techniques designed for non-reasoning LLMs largely fail to improve FC in LRMs. This paper identifies the critical need for advanced, LRM-specific prompt engineering strategies and proposes several promising future research directions, including mechanistic prompt design, adaptive prompting, and reinforcement learning for FC optimization, to bridge the observed faithfulness gap.
\end{abstract}

\section{Introduction}
The rapid advancements in Large Language Models (LLMs) have led to their deployment across diverse and often critical domains, from scientific discovery to legal analysis and medical diagnostics \cite{Gani_2026_Quantifying}. In such contexts, the reliability of model outputs, especially their communication of uncertainty, is paramount. Users frequently interpret the extended reasoning traces generated by Large Reasoning Models (LRMs) as evidence of deliberation, competence, and confidence, making the alignment between a model's internal epistemic state and its linguistic expression of certainty a crucial aspect of trustworthiness. This alignment is formally defined as Faithful Calibration (FC).

However, achieving FC in LRMs presents significant challenges. Unlike simpler LLM tasks, LRM outputs often involve long, multi-step reasoning processes with dynamic confidence shifts, inconsistent step structures, and complex conditional dependencies. Traditional methods for evaluating confidence, often focused on factual correctness or single-turn responses, do not adequately capture the nuanced nature of reasoning traces. Furthermore, prior work has shown that LLMs often exhibit a "faithfulness gap," conveying false information with decisive language (hallucinations) or overstating confidence \cite{Gani_2026_Quantifying}.

This paper delves into the findings of Gani et al. (2026), who introduce a novel framework to systematically quantify and analyze faithful confidence expression in LRMs. We highlight their key observations regarding the limitations of current LRMs and existing prompt engineering techniques, ultimately advocating for a dedicated research agenda in advanced prompt engineering strategies tailored to the unique characteristics of reasoning models.

\section{Quantifying Faithful Calibration in LRMs}

Gani et al. (2026) address the complexities of measuring FC in LRMs by proposing a multi-faceted framework. Their methodology centers on comparing an LRM's intrinsic confidence, $C(s_i)$, with its linguistically expressed decisiveness, $D(s_i)$, at each step $s_i$ of a reasoning trace $T = (s_1, s_2, \dots, s_n)$.

\subsection{Problem Formulation}
Step-level faithfulness $F(s_i)$ is defined as:
\begin{equation}
F(s_i) = 1 - |D(s_i) - C(s_i)|
\end{equation}
where $F(s_i) \in [0, 1]$, with 1 indicating perfect alignment. This step-level score is then aggregated to a trace-level score $FC(T)$ and further to a dataset-level metric, $cMFG^*$, a width-weighted conditional mean faithfulness gap, designed to provide stable and comparable estimates across models with varying confidence distributions.

\subsection{Intrinsic Confidence Estimation}
To capture different facets of intrinsic confidence, the framework employs three complementary estimators:
\begin{itemize}
    \item \textbf{Recurrent Confidence Chain (RCC):} This representation-based probe captures confidence not just locally but also propagates uncertainty from earlier steps to later ones. The recurrent confidence state $p_i$ for step $s_i$ is given by:
    \begin{equation}
    p_1 = q_1, \quad p_i = \delta q_i + (1 - \delta) p_{i-1}
    \end{equation}
    where $q_i$ is the local step confidence and $\delta$ is a smoothing parameter.
    \item \textbf{DeepConf:} A token log-probability estimator that reflects the stochasticity of the generation process, with peaked next-token distributions indicating higher certainty.
    \item \textbf{Sampling Consistency:} This method, adapted for reasoning traces, estimates confidence from the consistency of repeated samples of a step, conditioned on the original prompt and prior steps. A novel prefix-conditioned sampling approach is introduced to control for conditional dependencies and structural variation.
\end{itemize}

\subsection{Linguistic Decisiveness Estimation}
Linguistic decisiveness $D(s_i)$ is estimated using an LLM-as-a-Judge approach. A judge model (e.g., Gemini-2.5-Flash) is prompted with few-shot examples to score text for linguistic assertiveness on a scale of $[0, 1]$, reflecting human perception of hedging language.

\subsection{Prompt Interventions}
The study investigates whether specialized prompt interventions, previously shown to boost FC in non-reasoning LLMs, are effective for LRMs. These include a baseline prompt, a "perception" prompt to elicit faithful linguistic confidence, and a "metacognitive" system prompt paired with the perception prompt.

\section{Key Findings and Implications for Prompt Engineering}

The empirical study by Gani et al. (2026) across diverse models, datasets, and prompt strategies revealed several crucial insights:

\begin{enumerate}
    \item \textbf{Systemic FC Deficiencies in LRMs:} Current LRMs generally exhibit only moderate faithful calibration. They tend to maintain high linguistic decisiveness even on difficult tasks where their answers are frequently incorrect. This suggests a fundamental disconnect between internal uncertainty and external expression.
    \item \textbf{Reasoning Training Degrades Faithfulness:} Surprisingly, reasoning-specific training (e.g., supervised fine-tuning on Chain-of-Thought traces) was found to \textit{degrade} faithful calibration. Reasoning-tuned models produced longer traces with more hesitation and self-correction language, but these linguistic markers did not correspond to lower internal confidence, making the caution unfaithful to the model's true uncertainty.
    \item \textbf{Prompt Interventions Fail to Generalize:} Prompting methods that improved FC in non-reasoning LLMs largely failed to generalize to LRMs. While some prompts could improve accuracy, these gains did not translate to better faithful calibration. This indicates that merely instructing an LRM to express uncertainty does not reliably repair the relationship between its intrinsic confidence and linguistic decisiveness.
    \item \textbf{Estimator-Dependent Confidence Signals:} The study found that different intrinsic confidence estimators (RCC, DeepConf, Sampling Consistency) produced divergent faithfulness profiles on identical traces. This suggests that LRMs may not possess a single, universal "uncertainty signal" that is consistently verbalized, highlighting the complexity of capturing true internal confidence.
\end{enumerate}

These findings underscore a significant challenge for prompt engineering: simply crafting prompts that instruct LRMs to be more "cautious" or to "express uncertainty" is insufficient. The problem of FC in LRMs is deeply intertwined with their internal reasoning processes, training paradigms, and the very nature of their confidence signals.

\section{Future Research Directions in Prompt Engineering for FC}

Given the limitations of current prompt engineering techniques for improving FC in LRMs, new research directions are essential:

\begin{enumerate}
    \item \textbf{Mechanistic Prompt Design for FC:} Instead of relying on heuristic prompts, future work should focus on developing prompts informed by the mechanistic understanding of LRM internals. This could involve "meta-prompts" that guide the model on \textit{how} to leverage its hidden states or token probabilities to generate linguistically faithful uncertainty. For example, prompts could explicitly refer to the "strength of internal activation for alternative tokens" or "coherence of reasoning paths" to guide uncertainty expression.
    \item \textbf{Adaptive and Dynamic Prompting for Confidence Trajectories:} Current prompts are static. A more sophisticated approach would involve dynamic prompting systems that adapt during a reasoning trace. If real-time intrinsic confidence estimates (e.g., a low $p_i$ from RCC or high entropy in DeepConf) indicate high uncertainty at a particular step, the system could dynamically inject a sub-prompt encouraging more cautious language for the subsequent steps, thereby aligning $D(s_i)$ with $C(s_i)$ in real-time. This can be formalized as updating the prompt $P_t$ at time $t$ based on the confidence $C(s_{t-1})$ of the previous step.
    \item \textbf{Reinforcement Learning for FC-Aware Prompt Optimization:} Leveraging Reinforcement Learning from Human Feedback (RLHF) or AI Feedback (RLAIF) could enable the automatic discovery of prompts that optimize FC. By defining a reward function that incorporates the $FC(T)$ metric, models could be fine-tuned to generate responses that are both accurate and faithfully calibrated, moving beyond manual prompt engineering.
    \item \textbf{Contrastive Prompting for Uncertainty Introspection:} Design prompts that explicitly challenge the model to introspect on its uncertainty. For instance, asking the model to "provide the most likely answer and then explain why it might be wrong" or "generate an alternative reasoning path and its associated confidence" could force the LRM to engage in deeper self-assessment, potentially leading to better alignment between its internal state and expressed confidence.
    \item \textbf{Multimodal Prompting for Confidence Signals:} For Vision-Language Models (VLMs) or other multimodal LRMs, explore whether incorporating non-textual confidence signals (e.g., visual uncertainty cues, or even synthetic "confidence tokens" during input) into prompts can help bridge the FC gap. This could provide richer context for the model to ground its linguistic uncertainty.
\end{enumerate}

\section{Conclusion}
The work by Gani et al. (2026) highlights Faithful Calibration as a critical and under-examined alignment problem for Large Reasoning Models. Their comprehensive framework reveals that LRMs currently struggle to faithfully express their intrinsic confidence, and existing prompt engineering techniques are largely ineffective in addressing this gap. As LRMs are increasingly deployed in high-stakes environments, the need for reliable uncertainty communication becomes paramount. Future research must move beyond superficial prompt interventions and explore advanced, LRM-specific prompt engineering strategies that are deeply integrated with the models' internal mechanisms and dynamic reasoning processes. By focusing on mechanistic prompt design, adaptive prompting, and reinforcement learning for FC optimization, we can pave the way for more trustworthy and transparent AI systems.

\bibliographystyle{plain}
\begin{thebibliography}{1}

\bibitem{Gani_2026_Quantifying}
Areeb Gani, Asal Meskin, Gabrielle Kaili-May Liu, and Arman Cohan. Quantifying faithful confidence expression in large reasoning models. \textit{arXiv preprint arXiv:2606.03969}, 2026. \href{https://arxiv.org/pdf/2606.03969v1}{[PDF]}

\end{thebibliography}

\end{document}
